% Part: sets-functions-relations
% Chapter: induction
% Section: induction

\documentclass[../../include/open-logic-section]{subfiles}

\begin{document}

\olfileid{sfr}{ind}{int}

\olsection{جان پچھان}

\begin{explain}
استقرا ریاضی دی اک ثبوتی چال اے جیہڑی منطق وچ عام ورتی جاندی اے۔ ہو سکدا اے تسی ریاضی توں ایہدے نال واقف ہوو، جتھے قدرتی عدداں اُتے استقرا کردے او۔ اصل وچ ریاضی تے منطق دیاں شیواں دے کئی وکھ وکھ طرح دے سیٹ نیں جنہاں اُتے استقرا لگ سکدا اے۔

عام طور اُتے استقرا نال تسی اک کلی دعویٰ ثابت کر سکدے او، یعنی وکھا سکدے او کہ کسے خاص قسم دی ہر شے وچ کوئی خاصیت اے۔ خاص کر استقرا اکثر اوہدوں کم آؤندا اے جدوں «کلی شامل کرن» والا ثبوتی طریقہ بہت گنجل دار ہووے۔ استقرا صرف ایہو جیہیاں ریاضیاتی شیواں اُتے لگدا اے جیہڑیاں خاص طریقیاں نال بنیاں ہون۔ جے سیٹ دی ہر شے یا بنیادی اے یا بنیادی شیواں توں بنی اے، تاں اسیں اوس اُتے استقرا ورت سکدے آں۔ ہور ٹھیک طرح ایہ گل ایویں اے:
\end{explain}

\begin{defn}
سیٹ $S$ کسے فنکشن $f$ ہیٹھ \emph{بند} اے اوہدوں تے صرف اوہدوں جدوں ہر واری جے $x \in
S$، تاں $f(x) \in S$۔
\end{defn}

\begin{explain}
مثال لئی قدرتی عدد ($\mathbb{N}$) اگلا عدد دین والے فنکشن $f(x) = x+1$ ہیٹھ بند نیں۔ جے $x$ قدرتی عدد اے، تاں $x+1$ وی اے۔ پر قدرتی عدد پچھلا عدد دین والے فنکشن $f(x) = x-1$ ہیٹھ بند نہیں، کیونکہ $0$ قدرتی عدد اے، پر $-1$ نہیں۔
\end{explain}

\begin{defn}
آکھدے آں کہ خاصیت $P$، فنکشن $f(x)$ ہیٹھ \emph{قائم رہندی} اے جدوں کسے وی شے $a$ لئی جیہڑی $f$ دے دائرے وچ اے، $P(a)$ توں $P(f(a))$ نکلدا اے (جے $a$ وچ خاصیت $P$ اے، تاں $f(a)$ وچ وی اے)۔
\end{defn}

\begin{explain}
«جفت عدد ہونا» والی خاصیت فنکشن $f(x) = x+2$ ہیٹھ قائم رہندی اے، کیونکہ جے $x$ جفت اے، تاں $x+2$ وی جفت اے۔ جفت عدد ہون دی خاصیت اگلا عدد دین والے فنکشن ہیٹھ قائم نہیں رہندی، کیونکہ جے $x$ جفت اے، تاں $x+1$ طاق اے۔

استقرا قدرتی عدداں اُتے ایس لئی کم کردا اے کہ ہر قدرتی عدد تک 0 توں، اگلا عدد دین والا فنکشن محدود واریاں لا کے، پہنچیا جا سکدا اے۔ ایہ ہر اوس سیٹ دی خاص گل اے جیہدے اُتے اسیں استقرا کر سکدے آں۔
\end{explain}

\begin{defn}[$\Nat$ اُتے استقرا]
جے 0 وچ خاصیت $P$ اے تے جے $P$، اگلا عدد دین والے فنکشن ہیٹھ قائم رہندی اے، تاں ہر قدرتی عدد وچ خاصیت $P$ اے۔
\end{defn}

\begin{ex} پہلے $n$ قدرتی عدداں دا جوڑ $\frac{n(n+1)}{2}$ اے۔ \\

\textbf{مُڈھلی صورت۔} پہلے 0 قدرتی عدداں دا جوڑ $0=
\frac{0\cdot 1}{2}$ اے۔\\

\textbf{استقرائی قدم۔} منّ لو کہ خاصیت $k$ لئی ٹھیک اے۔ اسیں وکھاواں گے کہ ایہ $k+1$ لئی وی ٹھیک اے۔ دوجے لفظاں وچ، منّ لو کہ $P(k)$ ٹھیک اے، یعنی
\[ 0 + 1 + \ldots + k = \frac{k(k+1)}{2} \]
اسیں وکھاواں گے کہ $P(k+1)$ ٹھیک اے، یعنی
\[ 0 + 1 + \ldots + k + (k+1) = \frac{(k+1)(k+2)}{2} \]
ایہ مفروضہ ورت کے کہ $P(k)$ ٹھیک اے۔

\begin{align*}
0 + 1 + \ldots + k &= \frac{k(k+1)}{2} \tag{مفروضہ} \\
0 + 1 + \ldots + k + (k+1) &= \frac{k(k+1)}{2} + (k+1)
\tag{$k+1$ دوناں پاسیاں وچ پاؤ} \\
&= \frac{k(k+1) + 2(k+1)}{2} \\
&= \frac{2k + k + 2k +2}{2} \\
&=\frac{(k+1)(k+2)}{2}
\end{align*}
سو استقرائی قدم پورا ہو گیا، تے اسیں دعویٰ ثابت کر دِتّا اے۔
\end{ex}

\begin{explain}
!!{formula}s لئی بنیادی اجزا جوہری !!{formula}s نیں۔ ہور گنجل دار !!{formula}s گھٹ گنجل دار !!{formula}s نوں عمل کاراں نال جوڑ کے ملدے نیں۔ ہر عمل کار نوں !!{formula}s دے کسے فنکشن نال میل کھاندا سمجھ سکدے آں۔ اوہ فنکشن اک نواں !!{formula} دیندا اے جیہڑا ایہناں دوہاں نوں متعلقہ عمل کار نال جوڑدا اے۔ مثال لئی دو !!{formula}s $!!^a$ تے $!B$ نوں عطف وچ جوڑن والا فنکشن $ \mathcal E _\land (!!^a,
!B) = !!^a \land !B$ لکھیا جا سکدا اے۔ ایہناں نوں \emph{فارمولے بناؤن والے فنکشن} آکھ سکدے آں۔ !!{formula}s دا سیٹ ایہناں فنکشناں ہیٹھ بند اے۔ جدوں وی $!!^a$ تے $!B$، !!{formula}s نیں، تاں $\lnot !!^a$، $!!^a \land
!B$ وغیرہ وی نیں۔
\end{explain}

\begin{defn}[!!{formula}s اُتے استقرا]
جے ہر جوہری !!{formula} وچ خاصیت $P$ اے تے $P$، !!{formula} بناؤن والے فنکشناں ہیٹھ قائم رہندی اے، تاں ہر !!{formula} وچ خاصیت $P$ اے۔
\end{defn}

\begin{explain}
ایہ تعریف !!{formula}s اُتے استقرائی ثبوت دا اک «طریقہ» دَسّدی اے۔ جے سانوں وکھاؤن نوں آکھیا جاوے کہ ہر !!{formula} وچ خاصیت $P$ اے، تاں ہِٹھلے قدم چُکّو:\\

\textbf{مُڈھلی صورت۔} منّ لو $!!^a$ اک جوہری !!{formula} اے۔ [\ldots] سو $!!^a$ وچ خاصیت $P$ اے۔

\textbf{استقرائی قدم۔} منّ لو $!!^a$ تے $!B$، !!{formula}s نیں تے دوناں وچ خاصیت $P$ اے۔

صورت $\lnot$۔ [\ldots] سو $\lnot !!^a$ وچ خاصیت $P$ اے۔

صورت $\land$۔ [\ldots] سو $!!^a \land !B$ وچ خاصیت $P$ اے۔

صورت $\lor$۔ [\ldots] سو $!!^a \lor !B$ وچ خاصیت $P$ اے۔

صورت $\lif$۔ [\ldots] سو $!!^a \lif !B$ وچ خاصیت $P$ اے۔

صورت $\liff$۔ [\ldots] سو $!!^a \liff !B$ وچ خاصیت $P$ اے۔

صورت $\lforall[]$۔ [\ldots] سو $\lforall[x] !!^a$ وچ خاصیت $P$ اے۔

صورت $\lexists[]$۔ [\ldots] سو $\lexists[x] !!^a$ وچ خاصیت $P$ اے۔ \\

سو ہر !!{formula} وچ خاصیت $P$ اے۔
\end{explain}

\end{document}
